% Pour toute question sur l'utilisation de cette classe, écrivez à jerome.soucy(at)ulaval.ca
% Nécessite la classe Evaluation.sty dans le même dossier ou dans le PATH de votre système
% Nécessite le fichier logo_ul.pdf dans le même dossier un dans le PATH de votre système
\def\Ponderation{35}
\def\SigleNumero{STT-1000}
\def\NomCours{Probabilités et statistique}
\def\DateEvaluation{11 décembre 2022 de 9h00 à 11h50}
\def\Session{}
\def\NomEvaluation{Examen 2}

% Ne rien mettre entre les accolades si l'enseignant (ou l'enseignante) est aussi le coordonnateur (ou la coordonnatrice)
\def\Coordonnatrice{}
\def\Coordonnateur{}

% Ne rien mettre entre les accolades s'il y a plus d'une section
\def\Enseignant{}
\def\Enseignante{}

% Ne rien mettre entre les accolades s'il s'agit d'un cours à section unique
% Mettre le nom de chacun des enseignants et enseignantes entre les accolades
% Une case à cocher sera générée pour chacune des sections
\def\SectionA{Anne-Sophie Charest}
\def\SectionB{Julien Miron}
\def\SectionC{}
\def\SectionD{}

% Concerne seulement les travaux pratiques : écrire entre les accolades le nombre maximal de membres pour une équipe
\def\NombreLignes{}		

% Si vous souhaitez qu'apparaisse une déclaration d'honneur au bas de la première page, écrivez quelque chose entre les accolades
\def\EnLigne{}

\def\Directives{
\begin{enumerate}
\item Matériel autorisé:
\begin{itemize}
\item Calculatrice scientifique {\bf autorisée par la faculté des sciences et de génie};
\item Le feuillet de formules fourni avec l'examen. Celui-ci inclut la table de la loi normale.
\end{itemize}
\item Assurez-vous que cette \'evaluation comporte \numquestions ~questions réparties sur \numpages{}~pages, pour un total de \numpoints{}~points.\
\item \textbf{Veuillez ne pas débrocher votre examen.}
\item Veuillez déposer votre carte d'identité \textbf{admise} avec photo sur votre bureau.\
\item Vous devez\textbf{ justifier votre réponse} par une démarche claire, sauf pour les questions à choix de réponse.\
\item Pour les questions à choix de réponse, vous pouvez mettre un X, un crochet ($\checkmark$) ou remplir la case.
\item Le verso des feuilles peut être utilisé comme brouillon, \textbf{mais ne sera pas corrigé}. Attention à ne rien débrocher.
\end{enumerate}
}

% Veuillez indiquer le type d'évaluation entre crochets
% Examen : option à choisir pour un examen présentiel qui sera corrigé de manière classique. Un tableau des points sera affiché sur la page couverture.
% Gradescope : option à choisir pour un examen qui sera corrigé via Gradescope. Aucun tableau de pointage n'apparaîtra et il y a quelques différences au niveau de la mise en page
% TPS :
% TPF :
% Devoir :
\def\Gradescope{.} % Mode Gradescope activé (recto seulement). Mettre {} pour recto-verso.
\documentclass{Evaluation}
\usepackage{multicol,graphicx}
\begin{document}

\newpage
%%%%%%%%%%%%%%%% QUESTION %%%%%%%%%%%%%%%%
\begin{questions}

%CHAPITRE 9 - 20%

\question[8] Yasmine s'intéresse à une variable aléatoire $X$ suivant une loi Binomiale($n=70,p=0,3$). Elle souhaite calculer ${P( 18 < X < 35)}$, mais n'a pas accès à un ordinateur et n'a pas envie de faire tous les calculs à la main. En utilisant la méthode présentée en classe, calculez approximativement cette probabilité. 

\begin{solution}
On approxime la loi binomiale avec la loi normale, en utilisant la correction pour la continuité.\\

On s'intéresse à $X = \sum_{i=1}^n X_i$ où chaque $X_i$ suit une loi Bernoulli(p).\\
On sait donc que $E(X_i) = p =0.3$ et $V(X_i) = p(1-p)=0.3*0.7=0.21$. \\
Le théorème central limite nous indique que $X \approx N(\mu = 70*0,3=21 ;  \sigma^2= 70*0.21=14,7$). 

En utilisant la correction pour la continuité, on doit calculer $$P(17,5 \le X \le 35,5) \mbox{ où } X \sim N(21; 17,4).$$

On commence par standardiser: 
\begin{align*}
P(17,5 \le X \le 35,5) &= P\left(\frac{17,5-21}{\sqrt{17,4}} \le \frac{X-21}{\sqrt{17,4}} \le \frac{35,5-21}{\sqrt{17,4}}\right) \\
&= P( -0,839 \leq Z \leq 3,47) 
\end{align*}

On utilise maintenant la table de la loi normale:
\begin{align*}
P( -0,839 \leq Z \leq 3,47) & = P(Z \leq 3,47) - P(Z \leq -0,839) \\
& = 0,99974 - (1-0,79954) = 0,79928
\end{align*}



\end{solution}
\newpage
%Ici, je donne directement la distribution Binomiale pour éviter que certains se mélangent dès le départ. 
%3 pts bonne distribution; mu, sigma, normale
%2 pts standardiser
%2 pts lire table
%1pt correction continuité

\question[4] 
Jean-Claude suppose que le temps d'attente à l'hôpital suit une exponentielle de moyenne 240 minutes. 
Il veut construire un intervalle de confiance bilatéral à $95\%$ pour le vérifier.
Il procède donc comme suit. 
Il veut prend le temps d'attente de $16$ personnes choisies au hasard et connaître $\Bar{X}$ et $S^2$. 
Il obtient $\Bar{x}=250$ et $s^2=59'536$.
Comme il construit un intervalle de confiance bilatéral, il choisit $z_{\alpha/2}=1,96$, à partir de sa table de la loi normale.
Son intervalle de confiance à $95\%$ est donc $[130,44;369,56]$.
Comme celui-ci contient $240$, il se dit que cette valeur est plausible pour $\mu$. 
Après avoir vérifié les résultats de Jean-Claude, vous constatez qu'il y a une erreur dans ce qu'il a fait. 
Quelle est cette erreur?

\begin{solution}
   Comme $n=16<30$, on ne peut pas utiliser le TCL.
\end{solution}
\vspace{4cm}
\question[2]
%QCM - dans un exemple, choisir si paramètre ou estimateur; étonnamment pas si bien réussi l'an dernier
Le poids des tomates d’un agriculteur suit une loi normale de moyenne $300$ grammes
et de variance $100$ grammes$^2$. On prend des mesures sur 25 tomates choisies aléatoirement et on obtient une moyenne échantillonnale de 309 grammes et une variance échantillonnale de 118 grammes$^2$.
Est-ce que la valeur de 118 grammes$^2$ est un paramètre ou une statistique ?

\begin{oneparcheckboxes}
\choice Un paramètre
\choice Une statistique
\end{oneparcheckboxes}
\vspace{0.5cm}
\question[2]
Si $X \sim \chi^2_{24}$, alors $P(X<0)=P(X>0)=0.5$.

\begin{oneparcheckboxes}
\choice Vrai
\choice Faux
\end{oneparcheckboxes}
\vspace{0.5cm}
%volé de l'examen de reprise H22

\question[4]
Sachant que $\displaystyle \frac{(n-1)S^2}{\sigma^2}$ suit une loi du khi-carré à $n-1$ degrés de liberté, donnez $\text{Var}(S^2)$.

\begin{solution}
On peut calculer que 
\begin{align*}
Var(S^2) &= Var\left( \frac{\sigma^2}{(n-1)}\frac{(n-1)}{\sigma^2}S^2 \right)\\
    &= \frac{\sigma^4}{(n-1)^2} Var\left(\frac{(n-1)}{\sigma^2}S^2\right)  \\
    &= \frac{\sigma^4}{(n-1)^2}(2(n-1)) \\
    &= \frac{2\sigma^4}{(n-1)}
\end{align*}
\end{solution}

%CHAPITRE 10 - 40%

\newpage

\question[8]
%volé de reprise H22
On a construit un intervalle de confiance pour la proportion de personnes qui préfèrent le chocolat noir au chocolat au lait. On a utilisé un échantillon de $2400$ personnes, parmi lesquelles $960$ ont répondu qu’elles préféraient le chocolat noir. On a obtenu l’intervalle $[37, 83\%; 42, 17\%]$. Quel est le niveau de confiance de cet intervalle ?

\begin{solution}
La marge d'erreur de l'intervalle est 
$$ \frac{42,17-37,82}{2} = 2,175\%$$ 
On a donc 
\begin{align*}
0,02175 &= z_{\alpha/2}\sqrt{\frac{\hat{p}{(1-\hat{p}}}{n}} \\ 
&= z_{\alpha/2} \sqrt{\frac{960/2400{(1-960/2400)}}{2400}}\\
&= z_{\alpha/2} 0,01
\end{align*}
On trouve donc que $z_{\alpha/2} = 2,175$. La table nous indique que $\alpha/2 = 1-0.985 = 0.015$ et donc que $\alpha = 0.03$. Le niveau de confiance de l'intervalle est donc $0.97$. 
\end{solution}
\vspace{13cm}

\question[4]
%Volé de l'examen de reprise H22
On s'intéresse aux frais d'inscriptions des étudiants de 1er cycle à temps plein au Québec. 
Avec un échantillon aléatoire de $n=121$ étudiants de la population, on obtient l'intervalle de confiance à 95\% suivant: $[2\,550;2\,630]$. 
Cochez tous les énoncés qui sont vrais parmi les suivants. (Aucune justification requise)

\begin{checkboxes}
\choice Il y a 95\% des étudiants de 1er cycle à temps plein au Québec qui paient des frais d'inscription entre 2\,550 \$ et 2\,630 \$
\choice Si on pigeait un échantillon de 1\,000 étudiants pour construire un nouvel intervalle de confiance à 95\% sur la moyenne, celui-ci aurait une plus grande probabilité de contenir la vraie valeur de $\mu$ que celui obtenu à partir de 121 étudiants.
\choice Si on pigeait un autre échantillon de 121 étudiants pour construire un intervalle de confiance à 95\% sur la moyenne, l'intervalle obtenu aurait la même longueur que celui-ci.
\choice Si on pigeait plusieurs échantillons de 121 étudiants dans cette population, et qu'on construisait avec chacun d'eux l'intervalle de confiance à 95\% pour $\mu$, alors environ 95\% de ces intervalles contiendraient la vraie valeur moyenne des frais d'inscription des étudiants de 1er cycle à temps plein au Québec. 
\end{checkboxes}


\begin{solution}
Seul le dernier énoncé est vrai. 
\end{solution}

\newpage
\question[8]
On souhaite effectuer une étude afin de savoir ce que les gens préfèrent entre la fondue et la raclette. Pour cette étude, on supposera que les réponses des gens $X_1,\hdots,X_n$ sont indépendantes et identiquement distribuées selon une Bernoulli($p$), où $p$ représente la probabilité de choisir la fondue. Quelle doit être la taille d'échantillon minimale si on désire que notre intervalle de confiance à 90\% pour $p$ ait une longueur maximale de $0,04$?
\begin{solution}
    fff
\end{solution}


\newpage

\question
On suppose deux populations distribuées selon des lois normales de moyennes $\mu_1$ et $\mu_2$ et de variances $\sigma^2_1$ et $\sigma^2_2$ inconnues qui sont supposées égales. Vous avez les échantillons

\begin{align*}
X_{11}, \ldots, X_{1n_1} \overset{i.i.d.}{\sim} N(\mu_1, \sigma^2_1), & & X_{21}, \ldots, X_{2n_2} \overset{i.i.d.}{\sim} N(\mu_2, \sigma^2_2),
\end{align*}
où $n_1=8$ et $n_2=8$. Vous avez construit un intervalle de confiance à $95\%$ pour $\mu_1-\mu_2$ et obtenu $[0,1; 4,39]$. Lors de la présentation des résultats, vous réalisez que vous avez pris en note que $S^2_1=4,52$, mais vous ne trouvez plus la valeur de $S^2_2$. 
\begin{parts}
    \part[8] Quelle est la valeur de $S^2_2$?
    \vspace{14cm}
    \part[2] Que pouvez-vous dire au sujet de $\mu_1-\mu_2$?
\end{parts}



\begin{solution}
    
\end{solution}

\newpage
\question
Clara se déplace en autobus pour aller travailler et elle note minutieusement le temps nécessaire pour ses déplacements. Voici ses cinq derniers temps de déplacements, en minutes : 
$$ 22, 25, 24, 21, 33.$$
On peut supposer que ses temps de déplacements sont indépendants et suivent une loi normale de moyenne $\mu$ et de variance $\sigma^2$.

\begin{parts}
\part[6]
%1 pt moyenne, 2 pts s2, 2 pts quantiles, 1pt IC
Donnez un intervalle de confiance à 90\% pour la variance de ses temps de déplacement.
\begin{solution}
On doit d'abord calculer la variance échantillonnale. On trouve 
$$\bar{x} = \frac{1}{5}(22+25+24+21+33) = 25$$
On a ensuite 
\begin{align*}
s^2 &= \frac{1}{4}\left( (22-25)^2 + (25-25)^2 + (24-25)^2 + (21-25)^2 + (33-25)^2 \right)\\
&= \frac{1}{4}(9+0+1+16+64) \\
&= \frac{90}{4}\\
&= 22,5 \mbox{ minutes}^2    
\end{align*}  

Puisque $\alpha = 0.1$ on doit chercher dans la table $\chi^2_{0.05,4} = 9,49$ et $\chi^2_{0.95,4} = 0,71$. 
Notre intervalle est donc 
$$ \left[ \frac{4*22,5}{9,49} = 9,48, \frac{4*22,5}{0.71} =126,76 \right] $$
\end{solution}
\vspace{10cm}
\part[2] 
Est-ce une bonne idée d'utiliser le point milieu de cet intervalle comme estimateur ponctuel pour $\sigma^2$? Justifiez votre réponse.
%1pt pour non, 1 pt pour justification
\begin{solution}
Non. On sait que $s^2$ est non biaisé pour $\sigma^2$ et le centre de l'IC n'est pas $s^2$.
\end{solution}
\vspace{2.5cm}
\part[2]
Un collègue vous fait remarquer que la taille d'échantillon est très petite et que votre intervalle de confiance n'est probablement pas valide. Que lui répondez-vous?
\begin{solution}
Si les données suivent bien une loi normale, cet IC est valide peu importe la taille d'échantillon.
\end{solution}

\end{parts}

%CHAPITRE 11 - 40%
\newpage
%\textbf{Chapitre 11 - 40 points}


\question[]
Une entreprise veut connaître la proportion des gens qui utilise leur produit phare afin de tester les hypothèses $H_0: p = 0,40$ contre $H_1:p<0,4$. 
Elle sélectionne donc un échantillon de $n=36$ personnes et obtient une estimation de $\widehat{p}=0,385$. 
En utilisant $\alpha = 0,05$, la compagnie rejettera $H_0$ si
$$Z_0=\displaystyle \frac{\widehat{p}-p_0}{\sqrt{\displaystyle\frac{p_0(1-p_0)}{n}}}<-Z_{\alpha}=-1,64.$$


\begin{parts}
\part[8] Quelle est la puissance du test si la vraie valeur est $p=0,395$? 
\begin{solution}
    On a
    \begin{align*}
        1-\beta&=P(Z<-1,64|p_1=0,395)\\
        &=P\left(\left.\frac{\widehat{p}-0,4}{\sqrt{\displaystyle\frac{0,4*0,6}{9}}}<1,64 \right| p_1=0,38\right)\\
        &=P\left(\left.\widehat{p}<-1,64 \sqrt{\frac{0,4*0,6}{36}}+0,4\right|p_1=0,395\right)\\
        &= P\left(Z<\frac{-1,64 \sqrt{\displaystyle\frac{0,4*0,6}{36}}+0,005}{\sqrt{\displaystyle\frac{0,395*0,605}{36}}}\right)\\
        &= P\left(Z<-1,58214\right)\\
        &\approx 1-0,94295\\
        &\approx 0,05705
    \end{align*}
\end{solution}
\vspace{11cm}
\part[2] La puissance du test serait-elle plus grande ou plus petite si on avait un échantillon de $25$ observations au lieu de $10$ observations? 
(\textit{Aucune justification requise})
\begin{checkboxes}
\choice La puissance serait plus grande avec $n=25$.
\choice La puissance serait plus petite avec $n=25$.
\choice La puissance serait la même avec $n=25$ et $n=10$.
\choice On n'a pas assez d'information pour répondre à la question.
\end{checkboxes}

\part[2] La puissance du test serait-elle plus grande ou plus petite s'il avait décidé de faire un test bilatéral?
(\textit{Aucune justification requise})
\begin{checkboxes}
\choice La puissance serait plus grande avec un test bilatéral.
\choice La puissance serait plus petite avec un test bilatéral.
\choice La puissance serait la même avec un test bilatéral.
\choice On n'a pas assez d'information pour répondre à la question.
\end{checkboxes}

\end{parts}
\newpage
\question[4]
%QCM - Indiquer quelles hypothèses sont valides
%On peut modifier un peu.
Soit $\mu$ le temps moyen que prend une machine pour assembler des composants électroniques. Le responsable du contrôle de processus veut tester $H_0:\mu=20$ vs $H_1:\mu\neq 20$ avec un échantillon aléatoire de 100 temps lors d'une journée donnée. La moyenne de l'échantillon est $\Bar{x}=21,7$ secondes et que la valeur-p du test est de $0,0294$. Cochez tous les énoncés qui sont vrais parmi les suivants.
\begin{checkboxes}
\choice Nous sommes certains de rejeter $H_0$ pour n'importe quel seuil $\alpha>0$, car la valeur-p est inférieure à $0,05$.
\choice Nous rejetons $H_0$ au seuil $\alpha=0,01$
\correctchoice Nous rejetons $H_0$ au seuil $\alpha=0,1$
\correctchoice Si à la place nous faisions un test unilatéral avec $H_1>20$, nous rejetterions $H_0$ au seuil $\alpha=0,05$.
\choice Comme $\Bar{x}\neq20$, nous rejetons $H_0$ peu importe le seuil $\alpha$.
\end{checkboxes}

\vspace{0.5cm}

\question
On s'intéresse au salaire des résidents d'une grande tour à condo. On sait que la distribution de ces salaires est bimodale. Notre hypothèse est que le salaire moyen est de $80 000\$$ et on veut tester si c'est bien le cas au seuil $\alpha = 0.05$. Un échantillon de 50 résidents nous donne un salaire moyen de $70 000\$$, avec un écart-type de $28 000\$$. 
\begin{parts}

\part[4]
Expliquez, en mots, ce que représente une erreur de type II dans le contexte?

\begin{solution}
On fait une erreur de type II si on rejette $H_0$ alors que $H_1$ est vraie. Ici, on ferait une erreur de type II si on ne rejetait pas l'hypothèse nulle que le salaire moyen est $80 000\$$ alors que le salaire moyen est en effet différent de $80 0000\$$. 
\end{solution}
\vspace{5cm}

\part[8] 
Quel est le seuil observé (la valeur-p) pour ce test d'hypothèse? 
%3pts statistique de test
%2 pts loi normale et non t
%1 pt lire table
%2 pts bilatéral
\begin{solution}

La statistique de test est 

$$ Z_0 = \frac{\bar{x}-\mu_0}{s/\sqrt{n}} = \frac{70 000-80 000}{28 000/\sqrt{49}} =  -2.5 $$

Puisque le test est bilatéral, la valeur-p est 
$$ 2*P(Z > 2.5 ) = 2*( 1- 0.99379) = 0,01242$$
\end{solution}
\end{parts}

\newpage
\question[4]
Dans les situations suivantes, serait-il préférable d'utiliser un test de comparaison de moyennes pour des échantillons indépendants ou un test de comparaison de moyennes pour des échantillons appariés. 
\begin{parts}
    \part[1] On veut comparer le niveau en anglais des étudiants des cégeps de la région de Montréal à celui des étudiants des des cégeps de la région de Québec.
    \vspace{1.5cm}
    \part[1] On aimerait comparer la performance de deux de vélos de course en mesurant le temps qu'il faut pour monter le mont Bélair avec ces deux vélos.
    \vspace{1.5cm}
    \part[1] On veut comparer le temps moyen en secondes que prennent deux marques de four pour atteindre $200$ celsius.
    \vspace{1.5cm}
    \part[1] On veut connaître la fréquence cardiaque moyenne d'adultes dans la quarantaine avant et après un effort soutenu de 10 minutes.
\end{parts}
\vspace{1.5cm}
\question[2]
Pourquoi est-il indiqué sur la feuille de formules que la loi de la statistique de test pour une proportion est approximative?

\vspace{3cm}
\question[4]
Parmi les hypothèses suivantes, lesquelles sont valides dans la construction de tests comme on les présente dans le cours?

\begin{checkboxes}
\correctchoice $H_0 : \mu = 12$ vs $H_1 : \mu > 12$
\choice $H_0 : \mu \leq 12$ vs $H_1 : \mu > 12$
\choice $H_0 : \bar{x} = 4$ vs $H_1 : \bar{x} \ne 4$
\choice $H_0 : \sigma^2 < 1$ vs $H_1 : \sigma^2 > 1$
\correctchoice $H_0: \sigma = 1$ vs $H_1: \sigma \neq 1$
\end{checkboxes}
\newpage
\question[4]
Interpréter en contexte le résultat d'un test d'hypothèse. 

\end{questions}

\vspace{1in}
18,31.

\end{document} 
